\documentclass[10pt,a4paper]{amsart}
\usepackage[margin=19mm]{geometry}
\usepackage{amsmath,amssymb,mathtools}
\usepackage[scr=rsfs,cal=euler]{mathalfa}
\usepackage{xfrac}
\usepackage[unicode,hidelinks,bookmarks=false]{hyperref}
\newcommand{\ie}{i.e.\ }
\newcommand{\cf}{cf.\ }
\newcommand{\resp}{respectively\ }
\newcommand{\Subsection}{\subsection}
\providecommand{\nobreakdash}{\nobreak\hbox{-}}
\setlength{\emergencystretch}{2em}
\multlinegap0pt
\usepackage{bm}
\usepackage{bbm}
\usepackage{dsfont}
\usepackage{xspace}
\usepackage{cancel}
\usepackage{mathtools}
\usepackage{amsthm}
\makeatletter
\@ifundefined{theorem}{\newtheorem{theorem}{Theorem}}{}
\@ifundefined{lemma}{\newtheorem{lemma}[theorem]{Lemma}}{}
\makeatother
\newcommand{\EE}{\mathbb E}
\newcommand{\N}{\mathbb N}
\newcommand{\PP}{\mathbb{P}}
\title[The friendship paradox for trees]{The friendship paradox for trees}
\author{Rajat Subhra Hazra, Frank den Hollander, Nelly Litvak, Azadeh Parvaneh}
\date{Source: arXiv:2505.21774v2 (CC BY 4.0)}
\begin{document}
\maketitle

\noindent\emph{Source and licence.} The friendship paradox for trees. Version arXiv:2505.21774v2 (CC BY 4.0), distributed under the Creative Commons Attribution 4.0 International licence, as recorded in the arXiv licence metadata for this exact version and shown on the abstract page https://arxiv.org/abs/2505.21774v2. Full rights notice: licenses/arxiv-2505.21774-CC-BY-4.0.txt.\par\smallskip

\noindent\textit{Editorial scope.} Counted unit: the Lemma whose source label is \texttt{lem:randomindexSLLN} in the article (source0.tex lines 846--849, source label \texttt{lem:randomindexSLLN}), delivered together with the article's own complete original proof of it (source0.tex lines 851--1073, one \texttt{proof} environment of 223 lines). No mathematics is written, repaired or re-derived: statement and proof are the article's own text line for line. Required context retained uncounted: eq:KSthm* (lines 802--805). Every definition, display and label of those lines is retained unchanged. Omitted from this package and not redistributed: 1--801, 806--845, 850--850, 1074--1615. Nothing inside a retained excerpt is omitted or altered: every retained body is delivered exactly as the article's lines are written, so no omission receipt of any kind accompanies this package. Citations: the retained excerpts carry no citation command at all, so no bibliography entry of the article is transcribed and no reference is redistributed. Reference adaptation: none was needed and none was made; every reference inside the delivered unit resolves inside it, so no unresolved reference or citation is produced. Printed slips kept literal: no printed word, symbol, hypothesis or number of the counted statement or of its proof was altered. Layout and class: the article's own class is \texttt{article} with packages including babel, amsmath, amssymb, amsfonts, amsthm, mathrsfs, bbm, graphicx, xcolor, enumerate, geometry, marginnote, afterpage, fancyhdr; the wrapper is the lane's pinned \texttt{amsart} editorial wrapper at 10pt on A4, which restates the theorem-like environments the retained text needs and adds the article's own macro definitions that the retained text needs (\texttt{\textbackslash EE}, \texttt{\textbackslash N}, \texttt{\textbackslash PP}), with extra wrapper packages (bm, bbm, dsfont, xspace, cancel, mathtools). The delivered numbering is the unit's own. Assets: no retained span contains a figure environment, a graphics-inclusion command, a PDF-page inclusion, a table or a TikZ picture; no image file is copied into this package, the package declares no asset dependency and no image file is redistributed with it. Licence: version arXiv:2505.21774v2 of this article is distributed under the Creative Commons Attribution 4.0 International licence, as recorded in the arXiv licence metadata for this exact version and shown on the abstract page https://arxiv.org/abs/2505.21774v2. A full rights notice is delivered at licenses/arxiv-2505.21774-CC-BY-4.0.txt. Characters: every retained excerpt is pure ASCII, so the delivered engine, XeLaTeX with its default Latin Modern text font, reports no missing character.
\begin{equation}
\label{eq:KSthm*}
\frac{Z_m}{\mu^m} \stackrel{\PP-\mathrm{a.s.}}{\longrightarrow} W_\infty, \qquad m\to\infty,
\end{equation}

\begin{lemma}
\label{lem:randomindexSLLN}
$\frac{Z_m^\chi}{Z_{m-1}} \stackrel{\PP-\mathrm{a.s.}}{\longrightarrow} \mu f^\chi$, $m\to \infty$.
\end{lemma}

\begin{proof}
Conditional on $\mathcal{F}_{m-1}$, the random variables $Y_{m,1}^{\chi},\ldots,Y_{m,Z_{m-1}}^{\chi}$ are i.i.d., their common distribution does not depend on $m$, and 
\[
\EE[Y_{m,1}^{\chi}]=\mu f^{\chi} < \infty .
\]
For $a>0$, define
\[
\widetilde{Y}_{m,i}^{\chi} = Y_{m,i}^{\chi}\, \mathbbm{1}_{\{Y_{m,i}^{\chi}\leq a Z_{m-1}\}}, \qquad 1\leq i\leq Z_{m-1}.
\]
Then
\[
\frac{Z_m^{\chi}}{Z_{m-1}}-\mu f^{\chi} = A_m + B_m + D_m,
\]
where
\[
\begin{aligned}
A_m &= \frac1{Z_{m-1}} \sum_{i=1}^{Z_{m-1}} \Big( \widetilde{Y}_{m,i}^{\chi} - \EE[\widetilde{Y}_{m,i}^{\chi}\mid\mathcal{F}_{m-1}] \Big),\\
B_m &= \frac{1}{Z_{m-1}} \sum_{i=1}^{Z_{m-1}} Y_{m,i}^{\chi}\, \mathbbm{1}_{\{Y_{m,i}^{\chi}>aZ_{m-1}\}},\\
D_m &= \EE[\widetilde{Y}_{m,1}^{\chi}\mid\mathcal{F}_{m-1}] -\mu f^{\chi}.
\end{aligned}
\]
Hence, for any $\varepsilon > 0$,
\[
\begin{aligned}
&\PP\left( \Big| \frac{Z_m^\chi}{Z_{m-1}}-\mu f^\chi \Big| >\varepsilon \,\middle|\, \mathcal{F}_{m-1} \right)\\
&\qquad \leq \PP\left( |A_m|>\tfrac{\varepsilon}{3} \,\middle|\, \mathcal{F}_{m-1} \right) 
+ \PP\left( B_m >\tfrac{\varepsilon}{3} \,\middle|\, \mathcal{F}_{m-1} \right) 
+ \mathbbm{1}_{\{|D_m|>\tfrac{\varepsilon}{3}\}} \qquad \mathrm{a.s.}
\end{aligned}
\]
We treat the three terms separately.

\medskip\noindent
(1) By Chebyshev's inequality and conditional independence,
\[
\begin{aligned}
&\PP\left( |A_m|>\tfrac{\varepsilon}{3} \,\middle|\, \mathcal{F}_{m-1} \right)
\leq \frac{9}{\varepsilon^2}\, \EE\left[ A_m^2 \,\middle|\, \mathcal{F}_{m-1} \right]
= \frac{9}{\varepsilon^2 Z_{m-1}}\, \mathbb{V}{\rm ar}\left( \widetilde{Y}_{m,1}^{\chi} \mid \mathcal F_{m-1} \right)\\
&\leq \frac{9}{\varepsilon^2 Z_{m-1}}\, \EE\left[ (\widetilde{Y}_{m,1}^{\chi})^2 \,\middle|\, \mathcal{F}_{m-1} \right]
= \frac{9}{\varepsilon^2 Z_{m-1}}\, \EE\left[(Y_{m,1}^{\chi})^2\, \mathbbm{1}_{\{Y_{m,1}^{\chi}\leq aZ_{m-1}\}} 
\,\middle|\, \mathcal{F}_{m-1} \right] \qquad \mathrm{a.s.}
\end{aligned}
\]
For $k,m\in\N$, define
\[
\Omega_{m}(k) = \left\{\omega\in \Omega\colon\, k^{-1}\mu^{m-1}\leq Z_{m-1}(\omega)\leq k\mu^{m-1} \right\}.
\]
Then, since $\Omega_m(k)\in\mathcal F_{m-1}$,
\[
\begin{aligned}
\PP\left( \left\{|A_m|>\tfrac{\varepsilon}{3}\right\}\cap \Omega_{m}(k) \right)
&
=\EE\Big[ \mathbbm{1}_{\Omega_{m}}(k)\, \PP\left(|A_m|>\tfrac{\varepsilon}{3} \,\middle|\, \mathcal{F}_{m-1} \right) \Big]
\\&
\leq \frac{9}{\varepsilon^2}\, \EE\left[ \mathbbm{1}_{\Omega_m(k)}\frac{1}{Z_{m-1}}\, \EE\left[ (Y_{m,1}^{\chi})^2\, \mathbbm{1}_{\{Y_{m,1}^{\chi}\leq aZ_{m-1}\}}
\,\middle|\, \mathcal{F}_{m-1} \right] \right]
\\&
\leq \frac{9k}{\varepsilon^2}\,\mu^{-(m-1)}\, \EE\left[ (Y_{1,1}^{\chi})^2\, \mathbbm{1}_{\{Y_{1,1}^{\chi}\leq ak\mu^{m-1}\}}
\right].
\end{aligned}
\]
Hence
\[
\begin{aligned}
\sum_{m\in\N} \PP\left( \{|A_m|>\tfrac{\varepsilon}{3}\}\cap\Omega_m(k) \right)
&
\leq \frac{9k}{\varepsilon^2} \sum_{m\in\N} \mu^{-(m-1)}\, \EE\Big[ (Y_{1,1}^{\chi})^2\, \mathbbm{1}_{\{Y_{1,1}^{\chi}\le ak\mu^{m-1}\}} \Big]
\\&
= \frac{9k}{\varepsilon^2}\, \EE\Big[ (Y_{1,1}^{\chi})^2 \sum_{m\in\N} \mu^{-(m-1)}\,\mathbbm{1}_{\{Y_{1,1}^{\chi}\le ak\mu^{m-1}\}} \Big],
\end{aligned}
\]
where the last equality follows from Tonelli's theorem. Now, for every $y\in\N$, if
\[
y\leq ak\mu^{m-1},
\]
then
\[
m-1\geq \log_{\mu}\Big(\frac{y}{ak}\Big).
\]
Hence
\[
\begin{aligned}
\sum_{m\in\N} \PP\left( \{|A_m|>\tfrac{\varepsilon}{3}\}\cap\Omega_m(k) \right)
&
\leq \frac{9k}{\varepsilon^2}\,\EE\bigg[ (Y_{1,1}^{\chi})^2\,\mathbbm{1}_{\{Y_{1,1}^{\chi}>0\}}\, \sum_{m\geq \lceil\log_\mu(Y_{1,1}^{\chi}/(ak))\rceil+1} \mu^{-(m-1)}
\bigg]
\\&
\leq \frac{9k}{\varepsilon^2}\, \EE\bigg[ (Y_{1,1}^{\chi})^2\,\mathbbm{1}_{\{Y_{1,1}^{\chi}>0\}}\,  \mu^{-\log_{\mu}(Y_{1,1}^{\chi}/(ak))}\,(1-\mu^{-1})^{-1} \bigg]
\\&
\leq \frac{9ak^2}{\varepsilon^2(1-\mu^{-1})}\, \EE[Y_{1,1}^{\chi}] <\infty.
\end{aligned}
\]

\medskip\noindent
(2) Since the event $\{B_m>\tfrac{\varepsilon}{3}\}$ implies that at least one of the variables $Y_{m,1}^{\chi},\ldots,Y_{m,Z_{m-1}}^{\chi}$ is larger than $aZ_{m-1}$, we have, by the union bound,
\[
\PP\left( B_m>\tfrac{\varepsilon}{3} \,\middle|\, \mathcal{F}_{m-1}\right)
\leq Z_{m-1}\PP\left( Y_{m,1}^{\chi}>aZ_{m-1} \,\middle|\, \mathcal{F}_{m-1}\right)  \qquad \mathrm{a.s.}
\]
Hence, since $\Omega_m(k)\in\mathcal F_{m-1}$,
\[
\begin{aligned}
\PP\left(\{B_m>\tfrac{\varepsilon}{3}\}\cap\Omega_m(k)\right)
&
= \EE\left[ \mathbbm{1}_{\Omega_m(k)}\,\PP\left( B_m>\tfrac{\varepsilon}{3} \,\middle|\, \mathcal{F}_{m-1} \right) \right]
\\&
\leq \EE\left[ \mathbbm{1}_{\Omega_m(k)}\, Z_{m-1}\, \PP\left( Y_{m,1}^{\chi} > aZ_{m-1} \,\middle|\, \mathcal{F}_{m-1} \right) \right]
\\&
\leq k\mu^{m-1}\,\PP\left( Y_{1,1}^{\chi} > ak^{-1}\mu^{m-1} \right).
\end{aligned}
\]
Consequently,
\[
\begin{aligned}
\sum_{m\in\N} \PP\left( \{B_m>\tfrac{\varepsilon}{3}\}\cap\Omega_m(k) \right)
&
\leq k\, \sum_{m\in\N} \mu^{m-1}\, \PP\left( Y_{1,1}^{\chi} > ak^{-1}\mu^{m-1} \right)
\\&
= k\, \EE\bigg[ \sum_{m\in\N} \mu^{m-1}\, \mathbbm{1}_{\{Y_{1,1}^{\chi}>ak^{-1}\mu^{m-1}\}} \bigg],
\end{aligned}
\]
where the last equality follows from Tonelli's theorem. 
Now, for every $y\in\N$, if
\[
y > ak^{-1}\mu^{m-1},
\]
then
\[
m-1< \log_\mu\Big(\frac{yk}{a}\Big).
\]
Hence
\[
\begin{aligned}
\sum_{m\in\N}\PP\left(\{B_m>\tfrac{\varepsilon}{3}\}\cap\Omega_m(k)\right)
&
\leq
k\, \EE\bigg[\mathbbm{1}_{\{Y_{1,1}^{\chi}>0\}}\sum_{m\in\N:\, m-1<\log_{\mu}(kY_{1,1}^{\chi}/a)} \mu^{m-1} \bigg]
\\&
\leq k\, \EE\bigg[\mathbbm{1}_{\{Y_{1,1}^{\chi}>0\}}\, \frac{\mu^{\log_{\mu}(kY_{1,1}^{\chi}/a)+1}}{\mu-1}\bigg]
\\&
\leq \frac{\mu k^2}{a(\mu-1)}\,\EE[Y_{1,1}^{\chi}] < \infty.
\end{aligned}
\]
 
\medskip\noindent
(3) Note that
\[
D_m = -\EE\left[ Y^{\chi}_{m,1}\mathbbm{1}_{\{Y^{\chi}_{m,1}>aZ_{m-1}\}} \,\middle|\, \mathcal{F}_{m-1} \right].
\]
Hence
\[
\begin{aligned}
|D_m|\, \mathbbm{1}_{\Omega_m(k)}
&
= \EE\left[ Y_{m,1}^{\chi} \, \mathbbm{1}_{\{Y_{m,1}^{\chi}>aZ_{m-1}\}} \,\middle|\, \mathcal{F}_{m-1} \right]\, \mathbbm{1}_{\Omega_m(k)}
\\&
\leq \EE\left[ Y_{m,1}^{\chi}\, \mathbbm{1}_{\{Y_{m,1}^{\chi} > ak^{-1}\mu^{m-1}\}} \,\middle|\, \mathcal{F}_{m-1} \right]  
\\&
= \EE\left[ Y_{1,1}^{\chi}\, \mathbbm{1}_{\{Y_{1,1}^{\chi}>ak^{-1}\mu^{m-1}\}} \right]
\qquad \mathrm{a.s.}
\end{aligned}
\]
Since $\EE[Y_{1,1}^{\chi}]<\infty$, the dominated convergence theorem implies that
\[
\EE\left[ Y_{1,1}^{\chi}\, \mathbbm{1}_{\{Y_{1,1}^{\chi} > ak^{-1}\mu^{m-1}\}} \right] \to 0, \qquad m\to\infty.
\]
Hence, there exists an $m_0(\varepsilon,k)\in\N$ such that, for all $m\geq m_0(\varepsilon,k)$,
\[
|D_m| \, \mathbbm{1}_{\Omega_m(k)} < \tfrac{\varepsilon}{3} \qquad \mathrm{a.s.}
\]
Consequently,
\[
\PP\left( \{|D_m|>\tfrac{\varepsilon}{3}\}\cap\Omega_m(k) \right) = 0
\]
for all sufficiently large $m$, and therefore
\[
\sum_{m\in\N} \PP\left( \{|D_m|>\tfrac{\varepsilon}{3}\}\cap\Omega_m(k) \right) < \infty.
\]


\medskip 
Combining the estimates in (1)--(3), we obtain, for every fixed $k\in\N$,
\[
\sum_{m\in\N} \PP\left( \left\{ \left| \frac{Z_m^\chi}{Z_{m-1}} - \mu f^\chi \right|>\varepsilon \right\} \cap \Omega_m(k) \right) <\infty.
\]
For $k,M\in\N$, put
\[
\Omega_{k,M} = \bigcap_{m\geq M} \Omega_m(k).
\]
Since $\Omega_{k,M}\subseteq \Omega_m(k)$ for every $m\geq M$, it follows that
\[
\sum_{m\geq M} \PP\left( \left\{ \left| \frac{Z_m^\chi}{Z_{m-1}}-\mu f^{\chi} \right|>\varepsilon \right\} \cap \Omega_{k,M} \right) < \infty.
\]
Hence, if
\[
E_{\varepsilon}^{\chi} = \left\{\left\{ \left| \frac{Z_m^{\chi}}{Z_{m-1}}-\mu f^{\chi} \right| > \varepsilon \right\}
\,\,\, \mathrm{i.o.}\right\},
\]
then, by the Borel-Cantelli lemma,
\[
\PP\left(E_{\varepsilon}^{\chi} \cap \Omega_{k,M} \right)=0.
\]
By \eqref{eq:KSthm*},
\[
\PP\left( \bigcup_{k\in\N}\bigcup_{M\in\N}\Omega_{k,M} \right) = 1.
\]
Therefore
\[
\begin{aligned}
\PP\left(E_{\varepsilon}^{\chi}\right)
&=
\PP\left(E_{\varepsilon}^{\chi}\cap\left(\bigcup_{k\in\N}\bigcup_{M\in\N}\Omega_{k,M}\right)\right) 
\leq\sum_{k\in\N}\sum_{M\in\N}\PP\left(E_{\varepsilon}^{\chi}\cap \Omega_{k,M}\right)=0,
\end{aligned}
\]
and so
\[
\PP \left(\left\{ \left| \frac{Z_m^{\chi}}{Z_{m-1}}-\mu f^{\chi} \right| > \varepsilon \right\}
\,\,\, \mathrm{i.o.}\right) = 0.
\]
Since $\varepsilon>0$ was arbitrary, the claim in Lemma~\ref{lem:randomindexSLLN} follows.
\end{proof}

\end{document}
